% ACC 2027 -- v2 (math-complete draft)
% Changes from v1: all placeholder math expanded and made explicit; three
% propositions with proofs; real result tables; verified citations.
% Edit freely -- prose is deliberately short so it is easy to rewrite.

\documentclass[letterpaper,10pt,conference]{IEEEtran}

\IEEEoverridecommandlockouts
% \overrideIEEEmargins is provided by ieeeconf.cls (ACC author kit), not by
% IEEEtran.cls. Guarded so the draft compiles with either class; it becomes
% active automatically once ieeeconf.cls is on the TeX path.
\providecommand{\overrideIEEEmargins}{}
\overrideIEEEmargins

\usepackage{amsmath,amssymb,amsthm,bm}
\usepackage{booktabs,multirow,array}
\usepackage{graphicx}
\usepackage{xcolor}
\usepackage{cite}
\usepackage[hidelinks]{hyperref}

\graphicspath{{figures/}{figures/X2/}{figures/CF2/}{figures/all/}}

\newcommand{\R}{\mathbb{R}}
\newcommand{\T}{^{\mathsf T}}
\newcommand{\sat}{\operatorname{sat}}
\newcommand{\clip}{\operatorname{clip}}
\newcommand{\sgn}{\operatorname{sgn}}
\newcommand{\vecpart}{\operatorname{vec}}
\newtheorem{proposition}{Proposition}
\newtheorem{assumption}{Assumption}
\newtheorem{remark}{Remark}

\title{\LARGE \bf
Online PES-Adaptive Neuromorphic Membrane--Gate Residual Control for
Actuator-Constrained UAV Quaternion Attitude Tracking
}

\author{%
First A. Author$^{1}$, Second B. Author$^{2}$, and Third C. Author$^{3}$%
\thanks{$^{1}$Department/Institution. Email: author1@domain.com}%
\thanks{$^{2}$Department/Institution. Email: author2@domain.com}%
\thanks{$^{3}$Department/Institution. Email: author3@domain.com}%
}

\begin{document}

\maketitle
\thispagestyle{empty}
\pagestyle{empty}

\begin{abstract}
We present a baseline-preserving neuromorphic residual controller for
actuator-constrained quaternion attitude tracking. A classical quaternion
proportional--derivative (PD) law supplies the stabilizing body torque; a
population of paired continuous coupled membrane--gate (CMG) units, read
through a decoder adapted online by prescribed error sensitivity (PES),
supplies a bounded additive correction. The residual is bounded by
construction, so the closed loop never leaves the actuator-feasible set and
the baseline is recovered exactly when the residual authority is removed.
The controller is realized both directly and as a custom Nengo neuron type;
the two agree to $2.0\times10^{-6}$ relative mean error. On matched 45\,s
MuJoCo campaigns over five attitude missions and two rotor-actuated
platforms, mean quaternion-angle error improves by 17.3--60.6\% with zero
actuator saturation. A neuron-model benchmark isolates decoder adaptation,
not conductance adaptation, as the source of the gain, and a learning-rule
ablation shows that an annealed-eligibility variant matches instantaneous
PES on smooth demand while both diverge under a blend-form residual on
step commands, motivating the additive law. Physical Quanser
3DOF~Hover experiments reproduce the qualitative benefit at a deliberately
conservative residual authority ($+2.65\pm1.08\%$ over $n{=}3$ matched Roll
pairs) and delimit where simulation-to-hardware transfer does not hold.
\end{abstract}

\section{Introduction}
\label{sec:introduction}

Learning-enabled attitude control must satisfy two constraints that are
often in tension. Actuator commands must remain inside a feasible set that
is nonnegative and bounded, and the certified behavior of an existing
controller should not be discarded to obtain adaptation. Replacing a
stabilizing law with a learned policy forfeits both properties at once.

Residual architectures address this by adding a learned correction to a
fixed baseline \cite{johannink2019,silver2018,shi2021neural_aug}. Within
neuromorphic control, the Neural Engineering Framework (NEF) provides a
population-coded representation and the PES rule for online decoder
adaptation \cite{eliasmith2004neuraleng,bekolay2014nengo,voelker2015pes};
this combination has been demonstrated on adaptive manipulator control
\cite{dewolf2016arm} and on spiking flight controllers
\cite{stagsted2020,dupeyroux2021neuromorphic}. What remains open is a
residual formulation that is simultaneously (i) bounded so that actuator
feasibility is structural rather than empirical, (ii) adapted online with a
local rule, and (iii) validated on the same architecture in simulation and
on physical hardware.

\textbf{Contributions.}
\begin{enumerate}
\item A continuous, non-spiking coupled membrane--gate (CMG) unit with
paired directional current encoding, and a bounded additive residual law
for which actuator feasibility and exact baseline recovery are proved
(Propositions~\ref{prop:bounded}--\ref{prop:recovery}).
\item An online PES-adapted decoder with projection and leak, for which the
decoder iterate remains in a compact set and the decoded error contracts
(Proposition~\ref{prop:pes}).
\item A Nengo realization as a custom neuron type, numerically equivalent to
the direct implementation to $2.0\times10^{-6}$ relative mean error.
\item A matched five-mission, two-platform MuJoCo campaign with explicit
rotor allocation and saturation accounting, plus a neuron-model benchmark
that attributes the improvement to decoder adaptation rather than
conductance adaptation.
\item A learning-rule ablation (instantaneous PES versus annealed
eligibility) that is neutral on smooth demand and identifies a
saturation-driven failure of the blend-form residual on step commands,
which the additive law \eqref{eq:residual} removes structurally.
\item Physical 3DOF~Hover validation, reported with its negative cells and
its feasibility boundary.
\end{enumerate}

\textbf{Claim boundaries.} We do not claim superiority over all neuron
models on every mission, and no full-scale physical yaw360 was completed.

\section{Problem Formulation}
\label{sec:problem}

\subsection{Rigid-Body Attitude Dynamics}

Let $q\in\mathbb{S}^{3}$ be the body attitude quaternion, $\bm\omega\in\R^3$
the body angular velocity, $\bm J=\bm J\T\succ0$ the inertia, and
$\bm\tau\in\R^3$ the applied body torque:
\begin{align}
    \dot q &= \tfrac{1}{2}\,q\otimes[0,\bm\omega\T]\T,
    \label{eq:kin}\\
    \bm J\dot{\bm\omega} &= \bm\tau-\bm\omega\times(\bm J\bm\omega).
    \label{eq:dyn}
\end{align}

\subsection{Attitude Representation and the Double Cover}
\label{sec:doublecover}

Attitude evolves on $SO(3)$. We parameterize it by the unit quaternion
group $\mathbb{S}^3\subset\mathbb{H}$ under the Hamilton product
$\otimes$, with the rotation action on $\bm x\in\R^3$ given by
$[0,(\bm R(q)\bm x)\T]\T=q\otimes[0,\bm x\T]\T\otimes q^{-1}$ and
\begin{equation}
    \bm R(q)=\bm I_3+2q_0[\bm q_v]_\times+2[\bm q_v]_\times^2,
    \qquad q=[q_0,\bm q_v\T]\T,
    \label{eq:Rq}
\end{equation}
where $[\cdot]_\times$ is the skew operator. The map
$\bm R:\mathbb{S}^3\to SO(3)$ is a two-to-one covering,
$\bm R(q)=\bm R(-q)$, so every physical attitude has two antipodal
representations. Equivalently, in the even subalgebra
$\mathrm{Cl}^{+}_{3,0}$ the same rotation is a rotor
$\mathcal R=\exp(-\bm B\vartheta/2)$ acting by $\bm x\mapsto
\mathcal R\bm x\mathcal R^{\dagger}$, with $\bm B$ the unit bivector of the
rotation plane; $\mathcal R$ and $q$ are isomorphic and the same
sign ambiguity appears as $\mathcal R\sim-\mathcal R$. The exponential
parameterization
\begin{equation}
    q=\exp\!\Big(\tfrac{\vartheta}{2}[0,\bm n\T]\T\Big)
     =\Big[\cos\tfrac{\vartheta}{2},\;\bm n\T\sin\tfrac{\vartheta}{2}\Big]\T
    \label{eq:qexp}
\end{equation}
makes the covering explicit: $\vartheta$ and $\vartheta+2\pi$ give the same
$\bm R$ but opposite $q$. This is the structural reason a naive quaternion
feedback can drive the state the long way around the sphere.

\subsection{Quaternion Error and Tracking Objective}

For a desired attitude $q_d$ with body-frame rate $\bm\omega_d$, define
\begin{align}
    q_e &= q_d\otimes q^{-1},\qquad
    \bm e_q = \sgn(q_{e,0})\,\vecpart(q_e),
    \label{eq:qerr}\\
    \theta_q &= 2\cos^{-1}\!\big(\clip(|q_{e,0}|,0,1)\big)\in[0,\pi].
    \label{eq:theta}
\end{align}
By \eqref{eq:qexp}, $\bm e_q=\bm n\sin(\vartheta_e/2)$ and
$\theta_q=|\vartheta_e|$ is the eigenaxis angle: the geodesic distance on
$SO(3)$, i.e. the shortest rotation carrying $q$ to $q_d$
\cite{wie1985quaternion,wie1989eigenaxis}. The sign selection in
\eqref{eq:qerr} chooses the antipode with $q_{e,0}\ge0$, hence
$\theta_q\le\pi$, and prevents the closed loop from unwinding through a
full revolution to reach an attitude it already occupies
\cite{mayhew2013unwinding,mayhew2011acc}.

\begin{assumption}[Branch selection]
\label{ass:branch}
$\sgn(\cdot)$ in \eqref{eq:qerr} is evaluated with $\sgn(0)=+1$, so
$\bm e_q$ is well defined on all of $\mathbb{S}^3$.
\end{assumption}

Differentiating $q_e=q_d\otimes q^{-1}$ using \eqref{eq:kin} and
$\dot{q^{-1}}=-\tfrac12[0,\bm\omega\T]\T\otimes q^{-1}$ gives the error
kinematics
\begin{align}
    \dot q_e&=\tfrac{1}{2}\,q_e\otimes[0,\bm\omega_e\T]\T,
    \qquad
    \bm\omega_e=\bm\omega-\bm R(q_e)\T\bm\omega_d ,
    \label{eq:errkin}\\
    \dot q_{e,0}&=-\tfrac12\bm q_{e,v}\T\bm\omega_e,
    \qquad
    \dot{\bm q}_{e,v}=\tfrac12\big(q_{e,0}\bm I_3
        +[\bm q_{e,v}]_\times\big)\bm\omega_e ,
    \label{eq:errkin_split}
\end{align}
so the tracking objective is $(\bm e_q,\bm\omega_e)\to(\bm 0,\bm 0)$,
equivalently $\theta_q\to0$.

\begin{remark}
The discontinuity of $\sgn(q_{e,0})$ at $q_{e,0}=0$ means
\eqref{eq:qerr} yields asymptotic stability that is \emph{almost} global
but not robustly globally asymptotically stable: arbitrarily small
measurement noise near the $q_{e,0}=0$ set can induce chattering between
branches. Robust global tracking requires a hysteretic hybrid selection
rather than a static sign \cite{mayhew2011hybrid,mayhew2013unwinding}. We
retain the static form because the residual of Sec.~\ref{sec:cmg_nengo} is
bounded by Proposition~\ref{prop:bounded} and cannot itself excite the
switching set, but we do not claim robust global stability.
\end{remark}

Throughout, $\theta_q$ is the reported simulation metric; the physical
Hover summaries instead use the wrapped Euler-vector norm, and the two are
not interchangeable.

\subsection{Actuator-Constrained Control Objective}

Let $\bm u\in\R^m$ denote the $m$ bounded command channels with symmetric
limits $\bm\ell\in\R^m_{>0}$, and let
\begin{equation}
    \mathcal{U}=\{\bm u\in\R^m:\;|u_i|\le\ell_i,\;i=1,\dots,m\}.
    \label{eq:Uset}
\end{equation}
For a four-rotor platform the collective thrust $T$ and body torque
$\bm\tau$ map to rotor thrusts $\bm f\in\R^4_{\ge0}$ through the allocation
matrix $\bm B$ (Sec.~\ref{sec:baseline}). Feasibility therefore requires
both $\bm u\in\mathcal{U}$ and $\bm f=\bm B^{-1}[T,\bm\tau\T]\T$ inside the
per-rotor bounds $[0,f_{\max}]$. The design objective is to reduce
$\theta_q$ relative to the baseline while (i) keeping every command in
$\mathcal{U}$, (ii) incurring no rotor saturation, and (iii) bounding the
deviation from the baseline command.

\section{Baseline Quaternion Attitude Control}
\label{sec:baseline}

\subsection{Quaternion PD Law}

With per-axis natural frequencies $\bm\omega_n$ and damping $\zeta$,
\begin{align}
    \bm K_q &= 2\operatorname{diag}(\bm J)\odot\bm\omega_n^{2},
    \qquad
    \bm K_{\omega} = 2\zeta\operatorname{diag}(\bm J)\odot\bm\omega_n,
    \label{eq:gains}\\
    \bm\tau_0 &= \sat_{\bm\tau_{\max}}\!\big(
        -\bm K_q\odot\bm e_q-\bm K_{\omega}\odot\bm\omega\big).
    \label{eq:pd}
\end{align}
For the regulation case $\bm\omega_d\equiv\bm 0$ and in the absence of
saturation, the candidate
\begin{equation}
    V=\tfrac12\bm\omega\T\bm J\bm\omega
      +2k_q\big(1-|q_{e,0}|\big),\qquad k_q>0,
    \label{eq:lyap}
\end{equation}
gives $\dot V=-\bm\omega\T(\bm K_\omega\odot\bm\omega)\le0$ along
\eqref{eq:dyn}--\eqref{eq:pd}, which together with the shortest-path
selection yields almost-global asymptotic stability of $(\bm e_q,
\bm\omega)=(\bm 0,\bm 0)$. This is the classical quaternion feedback
regulator \cite{wie1985quaternion,hua2009,fresk2013quaternion}; when
$\bm K_q$ and $\bm K_\omega$ are scalar multiples of $\bm J$ the closed-loop
response is an eigenaxis rotation \cite{wie1989eigenaxis}. Gains
\eqref{eq:gains} place each axis at the prescribed $(\omega_{n,i},\zeta)$
under the diagonal approximation of $\bm J$. The saturation in
\eqref{eq:pd} is the mechanism through which actuator limits enter; the
residual of Sec.~\ref{sec:cmg_nengo} is designed so that this property is
retained exactly.

\subsection{Collective Thrust and Rotor Allocation}

For rotor positions $\bm p_i$ and reaction coefficients $\kappa_i$,
\begin{equation}
    \bm B=\begin{bmatrix}
      1&1&1&1\\
      p_{1,y}&p_{2,y}&p_{3,y}&p_{4,y}\\
      -p_{1,x}&-p_{2,x}&-p_{3,x}&-p_{4,x}\\
      \kappa_1&\kappa_2&\kappa_3&\kappa_4
    \end{bmatrix},
    \quad
    \bm f=\bm B^{-1}\begin{bmatrix}T\\\bm\tau\end{bmatrix}.
    \label{eq:mixer}
\end{equation}
Collective thrust remains entirely classical; the CMG residual augments
only $\bm\tau$. Consequently the neural path cannot alter the vertical
loop, and any rotor infeasibility is attributable to the torque channel
alone.

\section{Online PES-Adapted Membrane--Gate Residual Control}
\label{sec:cmg_nengo}

\subsection{Paired Directional Current Encoding}

Given the bounded baseline $\bm u_0$ and limits $\bm\ell$, let
$\bar{\bm u}=\bm u_0\oslash\bm\ell\in[-1,1]^m$ and
\begin{equation}
    \bm I^{\pm}=I_b\bm 1+I_g\max(\pm\bar{\bm u},\bm 0)\in\R^m ,
    \qquad
    \bm I=\begin{bmatrix}\bm I^{+}\\\bm I^{-}\end{bmatrix}\in\R^{2m}.
    \label{eq:encode}
\end{equation}
Each channel is thus represented by an antagonistic pair, so that sign
information survives a nonnegative activity readout. For UAV attitude,
$m=3$ and the population has $2m=6$ units.

\subsection{Coupled Membrane--Gate Dynamics}

For unit $j$ with membrane state $v_j$ and gate state $r_j$,
\begin{align}
    \dot r_j &= \frac{\sigma\!\big(\alpha_h[v_j-(v_{\mathrm{th}}
        -\gamma r_j)]\big)-r_j}{\tau_r},
    \label{eq:gate}\\
    \dot v_j &= -\frac{v_j-v_{\mathrm{rest}}}{\tau}
        +\frac{I_j}{\tau}-c_j\,(v_j-v_{\mathrm{reset}}),
    \label{eq:mem}\\
    c_j &= \big(g_0+\mathbf 1_{\mathrm{adapt}}\,g_1|I_j|\big)r_j^{2},
    \label{eq:cond}
\end{align}
with $\sigma(\cdot)$ the logistic function. The pair
\eqref{eq:gate}--\eqref{eq:cond} is mutually coupled: $r_j$ shapes the
membrane conductance $c_j$, while $v_j$ drives the gate through the
$\gamma$-shifted threshold. The model is continuous and non-spiking; no
reset event is emitted.

Both states are integrated by exponential Euler at step $h$, which is
unconditionally stable for \eqref{eq:mem} and is the form used in the
released implementation:
\begin{equation}
    v_j^{+}=v_je^{-\lambda_jh}
      +\frac{d_j}{\lambda_j}\big(1-e^{-\lambda_jh}\big),
    \label{eq:expeuler}
\end{equation}
with $\lambda_j=\tau^{-1}+c_j$ and
$d_j=\tau^{-1}(v_{\mathrm{rest}}+I_j)+c_jv_{\mathrm{reset}}$.

\subsection{Continuous Activity and Decoder}

\begin{align}
    \bm a_k &= \clip_{[0,1]}\!\big(G\,\bm r_k\big)\in[0,1]^{2m},
    \label{eq:act}\\
    \bm W_0 &= \begin{bmatrix}\bm I_m&-\bm I_m\end{bmatrix}\in\R^{m\times2m},
    \label{eq:W0}\\
    \bm s_k &= \clip_{[-1,1]}\!\big(\bm W_k\bm a_k\big)\in[-1,1]^m .
    \label{eq:decode}
\end{align}
Initialization \eqref{eq:W0} makes the decoder read the antagonistic
difference of each pair, so the untrained network is sign-consistent with
the baseline by construction.

\subsection{Online PES Decoder Adaptation}

With teaching gains $k_e,k_\omega^{*}\ge0$, learning rate $\eta>0$, leak
$\lambda_W\ge0$ and projection $\Pi$ onto
$[-W_{\max},W_{\max}]^{m\times2m}$,
\begin{align}
    \bm y_k^{*} &= \clip_{[-1,1]}\!\big(
        -k_e\bm e_{q,k}-k_\omega^{*}\bm\omega_k\big),
    \label{eq:target}\\
    \bm\varepsilon_k &= \bm y_k^{*}-\bm s_k,
    \label{eq:pes_err}\\
    \bm W_{k+1} &= \Pi\!\big(\bm W_k
        +\eta\,\bm\varepsilon_k\bm a_k\T-\lambda_W\bm W_k\big).
    \label{eq:pes}
\end{align}
The update is causal: $\bm s_k$ is decoded with $\bm W_k$, and
$\bm W_{k+1}$ affects only sample $k{+}1$. Equation \eqref{eq:pes} is the
NEF prescribed-error-sensitivity rule \cite{bekolay2014nengo,voelker2015pes}
augmented with projection and leak.

\subsection{Bounded Additive Residual}

\begin{align}
    \bm u_k^{\mathrm{neural}} &= \bm s_k\odot\bm\ell,
    \label{eq:uneural}\\
    \bm u_k &= \sat_{\bm\ell}\!\big(
        \bm u_{0,k}+\beta\,\bm u_k^{\mathrm{neural}}\big),
    \qquad \beta\in[0,1],
    \label{eq:residual}\\
    \bm\delta_k &= \bm u_k-\bm u_{0,k}.
    \label{eq:delta}
\end{align}
Equation \eqref{eq:residual} is baseline-preserving and is not the
historical blend $(1-\beta)\bm u_0+\beta\bm u^{\mathrm{neural}}$, which
attenuates the baseline and can produce apparent gains by detuning alone.

\subsection{Structural Guarantees}

\begin{proposition}[Feasibility and bounded authority]
\label{prop:bounded}
For every $k$ and any $\bm W_k$, $\bm u_k\in\mathcal U$ and
$\|\bm\delta_k\|_\infty\le\beta\max_i\ell_i$.
\end{proposition}
\begin{proof}
$\bm u_k\in\mathcal U$ is immediate from the outer $\sat_{\bm\ell}$ in
\eqref{eq:residual}. From \eqref{eq:decode}, $\|\bm s_k\|_\infty\le1$, so
by \eqref{eq:uneural} $|\beta u^{\mathrm{neural}}_{k,i}|\le\beta\ell_i$.
Since $\sat_{\bm\ell}$ is $1$-Lipschitz and $\bm u_{0,k}\in\mathcal U$
implies $\sat_{\bm\ell}(\bm u_{0,k})=\bm u_{0,k}$,
$|\delta_{k,i}|=|\sat_{\ell_i}(u_{0,k,i}+\beta u^{\mathrm{neural}}_{k,i})
-\sat_{\ell_i}(u_{0,k,i})|\le\beta\ell_i$.
\end{proof}

Proposition~\ref{prop:bounded} is the reason actuator feasibility is
structural: the clip in \eqref{eq:decode} bounds the decoder output
\emph{before} it is scaled, so no decoder state, including a diverged one,
can command outside $\mathcal U$ or perturb the baseline by more than
$\beta\bm\ell$.

\begin{proposition}[Exact baseline recovery]
\label{prop:recovery}
If $\beta=0$, or if $\bm W_k=\bm 0$, then $\bm u_k=\bm u_{0,k}$ for all
$k$.
\end{proposition}
\begin{proof}
Either condition gives $\beta\bm u^{\mathrm{neural}}_k=\bm 0$ in
\eqref{eq:residual}; then $\bm u_k=\sat_{\bm\ell}(\bm u_{0,k})=\bm u_{0,k}$.
\end{proof}

\begin{proposition}[Bounded decoder and contraction]
\label{prop:pes}
The iterate \eqref{eq:pes} satisfies
$\bm W_k\in[-W_{\max},W_{\max}]^{m\times2m}$ for all $k$. Moreover, in the
unsaturated, unprojected regime with $\lambda_W=0$ and $\bm a_k=\bm a$
fixed, the decoded error obeys
\begin{equation}
    \bm\varepsilon_{k+1}
    =\big(1-\eta\|\bm a\|_2^{2}\big)\,\bm\varepsilon_k ,
    \label{eq:contract}
\end{equation}
so $\|\bm\varepsilon_k\|\to0$ monotonically iff
$0<\eta\|\bm a\|_2^{2}<2$.
\end{proposition}
\begin{proof}
Boundedness is immediate since $\Pi$ is applied at every step. For
\eqref{eq:contract}, with $\bm y^{*}$ fixed and no clipping,
$\bm s_{k+1}-\bm s_k=(\bm W_{k+1}-\bm W_k)\bm a
=\eta\bm\varepsilon_k\bm a\T\bm a=\eta\|\bm a\|_2^2\bm\varepsilon_k$.
Subtracting from $\bm y^{*}$ gives
$\bm\varepsilon_{k+1}=\bm\varepsilon_k-\eta\|\bm a\|_2^2\bm\varepsilon_k$.
\end{proof}

Equation \eqref{eq:contract} recovers the known PES convergence condition
\cite{voelker2015pes} and fixes the admissible learning-rate range from the
measured activity norm. Since $\bm a\in[0,1]^{2m}$, a sufficient condition
is $\eta<2/(2m)$.

\subsection{Adaptive Conductance Versus Adaptive Decoder}

Setting $\mathbf 1_{\mathrm{adapt}}=1$ in \eqref{eq:cond} leaves the
closed-loop trajectories numerically indistinguishable from
$\mathbf 1_{\mathrm{adapt}}=0$ in the amplified-gate regime, because
$c_j\propto r_j^{2}$ with $r_j=O(10^{-6})$ before amplification, so the
$g_1|I_j|$ term multiplies a vanishing quantity. The improvement reported
in Sec.~\ref{sec:mujoco_results} is therefore attributable to the online
decoder adaptation \eqref{eq:pes}, not to conductance adaptation.

\subsection{Nengo Realization and Numerical Equivalence}

The unit \eqref{eq:gate}--\eqref{eq:expeuler} is implemented as a custom
\texttt{nengo.neurons.NeuronType} with $\bm I$ injected directly into the
neuron input, the built-in \texttt{nengo.PES} rule supplying
\eqref{eq:pes}, and projection/leak applied after each simulator step. The
Nengo learning rate is converted analytically from $\eta$ by
$\eta_{\mathrm{nengo}}=\eta\,(2m)/h$ to match the internal $h/n$ scaling.
Against the direct implementation the isolated relative command error is
$1.05\times10^{-7}$, and the worst end-to-end mission mean-error difference
is $2.013\times10^{-6}$ relative.

\subsection{Timing and Telemetry}

MuJoCo runs use $h=0.5$\,ms for plant and neural update ($2$\,kHz). The
Hover hardware loop runs the Q8 interface at $200$\,Hz with ten $0.5$\,ms
Nengo substeps per hardware sample, preserving the $2$\,kHz neural rate.
Logged per step: $\bm v$, $\bm r$, $\bm a$, $\bm I$, $c_j$, the effective
threshold, $\bm s_k$, $\bm\varepsilon_k$, $\bm W_k$ and $\Delta\bm W_k$.

\section{Platform and Experimental Protocol}
\label{sec:platform_method}

\textbf{Platforms.} (i) Skydio~X2 MuJoCo model \cite{todorov2012mujoco} with
four named rotor actuators and allocation \eqref{eq:mixer}; (ii) an explicit
four-rotor Crazyflie design variant at $1.5\times$ nominal rotor authority.
The second is a controller/actuator co-design study and is labeled as such
throughout; it is not stock CF2 hardware. (iii) The Quanser 3DOF~Hover, with
encoder-tared attitude, filtered body rates, torque-to-$\{F,B,L,R\}$ voltage
allocation, amplifier gain $3$, command range $[0,24]$\,V, and residual
magnitude/slew caps.

\textbf{Missions.} \emph{Hover} ($q_d\equiv q_0$); \emph{Roll}
($\phi_d=0.25\sin0.7t$); \emph{Full3} (three-axis sinusoid, $0.5$\,rad);
\emph{Hard45} (mixed-frequency three-axis, $0.45$--$0.65$\,rad);
\emph{Hardstep} (piecewise-constant setpoints); \emph{Yaw360} (smoothstep
$0\!\to\!2\pi$ yaw with $\phi_d=\theta_d=0$).

\textbf{Matched protocol.} Baseline and CMG runs share reference, duration
($45$\,s), scale, allocator, initial condition, safety limits and logging.
Reported simulation metric is the mean of \eqref{eq:theta}; hardware
summaries use the wrapped Euler-vector norm. Every table states $n$.

\section{MuJoCo Results}
\label{sec:mujoco_results}

\begin{table}[!t]
\caption{Matched 45\,s missions, mean quaternion-angle error
$\bar\theta_q$ [rad], $h=0.5$\,ms, $n=1$ per cell. Hover is omitted: the
equilibrium is already exact and the ratio is undefined.}
\label{tab:mujoco}
\centering
\small
\begin{tabular}{llrrr}
\toprule
Platform & Mission & PD & CMG--PES & Improvement \\
\midrule
\multirow{5}{*}{Skydio X2}
 & Roll     & 0.050565 & 0.041785 & 17.36\% \\
 & Full3    & 0.303275 & 0.131748 & 56.56\% \\
 & Hard45   & 0.294082 & 0.138007 & 53.07\% \\
 & Hardstep & 0.454528 & 0.359355 & 20.94\% \\
 & Yaw360   & 0.075863 & 0.029902 & \textbf{60.58\%} \\
\midrule
\multirow{5}{*}{\shortstack[l]{CF2 rotor design\\($1.5\times$ authority)}}
 & Roll     & 0.051337 & 0.042439 & 17.33\% \\
 & Full3    & 0.288383 & 0.126476 & 56.14\% \\
 & Hard45   & 0.285449 & 0.132319 & 53.65\% \\
 & Hardstep & 0.426182 & 0.347886 & 18.37\% \\
 & Yaw360   & 0.075812 & 0.029870 & \textbf{60.60\%} \\
\bottomrule
\end{tabular}
\end{table}

Table~\ref{tab:mujoco} reports the primary campaign. Improvement is
consistent across both platforms and is largest on the missions with
sustained, smoothly varying demand (Yaw360, Full3, Hard45) and smallest on
the piecewise-constant Hardstep, where the tracking error is dominated by
transients that a bounded residual cannot pre-empt. No rotor saturation
occurred in any X2 cell.

\begin{figure}[!t]
\centering
\includegraphics[width=\linewidth]{X2/Hard45/x2_hard45_tracking.png}
\caption{X2, Hard45: matched PD and CMG--PES attitude tracking.}
\label{fig:x2_hard45_tracking}
\end{figure}

\begin{figure}[!t]
\centering
\includegraphics[width=\linewidth]{X2/Yaw360/x2_yaw360_neuron_activity.png}
\caption{X2, Yaw360 (largest improvement, 60.58\%): membrane state $\bm v$,
gate state $\bm r$, decoded activity $\bm a$, and decoder adaptation.
Antagonistic pairs are visible as complementary activation.}
\label{fig:x2_yaw360_neurons}
\end{figure}

\textbf{Ablation.} Replacing \eqref{eq:pes} by a fixed decoder
($\eta=0$) removes essentially all of the improvement, while toggling
$\mathbf 1_{\mathrm{adapt}}$ in \eqref{eq:cond} changes the result
negligibly. This isolates decoder adaptation as the operative mechanism.

\textbf{Six-DOF extension.} Extending the residual to $m=6$ (three force,
three torque channels) inside a position$\to$attitude cascade preserves the
structure of Sec.~\ref{sec:cmg_nengo}. On X2 the demanding coupled mission
improves position error by $76.1\%$ and attitude error by $65.8\%$; on the
lighter CF2 variant the same configuration yields no net benefit, which we
report as a platform-dependent limit rather than a tuning artifact.

\section{Learning-Rule Ablation: Annealed Eligibility}
\label{sec:eprop}

\subsection{Rule and Scope}

PES \eqref{eq:pes} updates the decoder from the instantaneous activity.
A natural alternative accumulates an eligibility trace and anneals both the
trace bandwidth and the step size, in the spirit of eligibility-propagation
methods for recurrent spiking networks \cite{bellec2020eprop}. With
$s=\clip(t/T,0,1)$ over a run of length $T$,
\begin{align}
    \rho(s) &= \rho_0-(\rho_0-\rho_1)\,s,
    &&\rho_0=0.98,\;\rho_1=0.50,
    \label{eq:decay}\\
    \bm z_{k} &= \rho(s_k)\,\bm z_{k-1}+\big(1-\rho(s_k)\big)\bm a_k,
    &&\bm z_{-1}=\bm 0,
    \label{eq:elig}\\
    \eta(s) &= \eta_0\,(1-0.8\,s),
    \label{eq:anneal}\\
    \bm W_{k+1} &= \Pi\!\big(\bm W_k+\eta(s_k)\,
        \bm\varepsilon_k\bm z_k\T-\lambda_W\bm W_k\big).
    \label{eq:eprop}
\end{align}
Equation \eqref{eq:eprop} reduces to \eqref{eq:pes} for $\rho\equiv0$ and
constant $\eta$. The schedule is coarse-to-fine: early updates integrate
activity over a long window at a large step, late updates approach the
instantaneous rule at $20\%$ of the initial step.

\emph{Scope.} These runs were produced by the precursor implementation of
the architecture, which differs from Sec.~\ref{sec:cmg_nengo} in two
respects that must be stated before the numbers are read. First, it applies
the \emph{blend} law $\bm\tau=(1-\beta)\bm\tau_0+\beta\bm\tau^{\mathrm{neural}}$
rather than the baseline-preserving additive law \eqref{eq:residual}, so
Proposition~\ref{prop:recovery} does not apply to it and part of any
fixed-decoder gain is baseline attenuation rather than neural correction.
Second, the actuation path is an ideal body-wrench injection rather than the
rotor allocation \eqref{eq:mixer}. Table~\ref{tab:eprop} is therefore
\emph{not} comparable cell-by-cell with Table~\ref{tab:mujoco}; the valid
comparison it supports is between the two learning rules under identical
conditions.

\subsection{Results}

\begin{table}[!t]
\caption{Learning-rule ablation under the precursor blend architecture
(ideal wrench actuation). Mean quaternion-angle error $\bar\theta_q$ [rad],
45\,s, $h=0.5$\,ms, $n=1$. Percentages are relative to PD. Not comparable
to Table~\ref{tab:mujoco} (different residual law and actuation path).}
\label{tab:eprop}
\centering
\small
\begin{tabular}{llrrrr}
\toprule
Plat. & Mission & PD & PES & e-prop & $\Delta$PES / $\Delta$e-prop \\
\midrule
\multirow{5}{*}{X2}
 & Roll     & 0.05056 & 0.01899 & 0.01518 & $+62.5\%$ / $\bm{+70.0\%}$ \\
 & Full3    & 0.30327 & 0.09956 & 0.09612 & $+67.2\%$ / $\bm{+68.3\%}$ \\
 & Hard45   & 0.29408 & 0.12205 & 0.12753 & $\bm{+58.5\%}$ / $+56.6\%$ \\
 & Hardstep & 0.27942 & 1.69357 & 1.87203 & $-506\%$ / $-570\%$ \\
 & Yaw360   & 0.07586 & 0.02706 & 0.02098 & $+64.3\%$ / $\bm{+72.3\%}$ \\
\midrule
\multirow{5}{*}{CF2}
 & Roll     & 0.05134 & 0.00258 & 0.00254 & $+95.0\%$ / $\bm{+95.1\%}$ \\
 & Full3    & 0.28838 & 0.02663 & 0.02608 & $+90.8\%$ / $\bm{+91.0\%}$ \\
 & Hard45   & 0.28545 & 0.03555 & 0.03544 & $+87.6\%$ / $\bm{+87.6\%}$ \\
 & Hardstep & 0.66187 & 1.75781 & 2.07947 & $-166\%$ / $-214\%$ \\
 & Yaw360   & 0.07581 & 0.00061 & 0.00065 & $\bm{+99.2\%}$ / $+99.1\%$ \\
\bottomrule
\end{tabular}
\end{table}

\begin{figure}[!t]
\centering
\includegraphics[width=\linewidth]{eprop/eprop_ablation_summary.png}
\caption{Learning-rule ablation of Table~\ref{tab:eprop}. The two rules are
within a few points on every smooth-demand mission; both diverge on
Hardstep, where the allocation saturates.}
\label{fig:eprop_summary}
\end{figure}

Table~\ref{tab:eprop} and Fig.~\ref{fig:eprop_summary} support three
observations. First, the two rules are
close on every smooth-demand mission: the annealed trace wins six of ten
cells, the instantaneous rule wins four, and the median absolute difference
is under two percentage points. The decoder becoming adaptive at all is what
matters; which local rule performs the adaptation is second order.

Second, both rules fail on Hardstep, and the annealed rule fails harder
($-570\%$ versus $-506\%$ on X2). The mechanism is visible in the actuator
channel: the saturated fraction rises from $0.089$ under PD to $0.453$ and
$0.553$ under PES and e-prop respectively. A proportional teaching target
\eqref{eq:target} on a piecewise-constant reference learns additional
effective gain, which overshoots at each step transition, saturates the
allocation, and does not recover. The longer eligibility window integrates
that transient over a wider support, which is why annealing makes it worse
rather than better.

Third, the reproduction is exact. Re-running the frozen configuration
recovers the archived composite scores to six decimal places
($0.287391$, $0.038258$, $0.038215$ for PD, PES, and e-prop on CF2/Hard45),
which fixes the numerical provenance of this comparison.

\subsection{Consequence for the Proposed Law}

The Hardstep failure is a property of the blend architecture, not of the
teaching signal alone: under $(1-\beta)\bm\tau_0+\beta\bm\tau^{\mathrm{neural}}$
a diverging decoder simultaneously attenuates the stabilizing baseline, so
there is no lower bound on retained baseline authority. The additive law
\eqref{eq:residual} removes this failure mode by construction
(Prop.~\ref{prop:bounded}--\ref{prop:recovery}): the baseline term is never
scaled, and the worst-case deviation is $\beta\bm\ell$ regardless of decoder
state. This ablation is therefore the empirical motivation for adopting
\eqref{eq:residual} in the proposed controller.

\section{Neuron-Model Benchmark}
\label{sec:neuron_benchmark}

Under matched paired-current encoding \eqref{eq:encode} and an identical
PES decoder, we compare CMG against LIF, adaptive LIF, Izhikevich
\cite{izhikevich2003} and LIFRate. Conventional models are read through a
$5$\,ms low-pass of their native output, calibrated so that $\bm W_0$,
$\beta$ and the residual retain identical meaning.

Individual cells can favor Izhikevich--PES, and we do not claim uniform
superiority. The consistent difference is in the cost of the correction:
CMG--PES attains its improvement with lower normalized residual RMS and
less decoder movement. In the 6-DOF coupled mission the spiking models
destabilize even at $1/50$ of the CMG residual authority, whereas the
non-spiking LIFRate remains stable, which localizes the failure to spike
discretization rather than to the tuning curve.

\section{Physical 3DOF Hover Results}
\label{sec:hover_results}

The hardware loop executes the Nengo network directly at $200$\,Hz with a
$2$\,kHz neural substep and full telemetry; the browser dashboard is
excluded from the real-time path because contention with the control thread
caused deadline violations.

\begin{table}[!t]
\caption{Physical 3DOF~Hover, matched pairs. Mean wrapped Euler-vector
error [deg]. Roll is $n=3$; remaining rows are $n=1$.}
\label{tab:hover}
\centering
\small
\begin{tabular}{lrrr}
\toprule
Mission & PD & CMG--PES & Improvement \\
\midrule
Roll ($n{=}3$) & $5.4405\pm0.0648$ & $5.2960\pm0.0428$
  & $+2.65\pm1.08\%$ \\
Hover   & 0.116  & 0.046  & $+60.38\%$\,$^{\dagger}$ \\
Full3   & 11.282 & 11.644 & $-3.21\%$ \\
Hard45  & 9.743  & 9.689  & $+0.55\%$ \\
Hardstep& 13.626 & 13.497 & $+0.95\%$ \\
Yaw (reduced, $0.3\times$) & 2.731 & 2.960 & $-8.38\%$ \\
\bottomrule
\multicolumn{4}{l}{\footnotesize $^{\dagger}$ near-zero baseline; the ratio
is not a meaningful effect size.}
\end{tabular}
\end{table}

The repeated Roll experiment is the primary physical result and is the only
row with replication. Two mission cells are negative. The residual cap and
slew limit on hardware are deliberately far below the simulation authority,
which bounds the attainable improvement by Proposition~\ref{prop:bounded}
and largely explains the magnitude gap with Table~\ref{tab:mujoco}.

\textbf{Yaw feasibility boundary.} Two full-scale PD attempts aborted on the
roll safety limit at approximately $\pm20.127^\circ$, at the same attitude
independent of sweep duration; only a reduced-range profile terminating near
$108^\circ$ completed. We report this as an unresolved feasibility boundary
and offer cable wind-up at the rotating joint as a hypothesis only.

\section{Discussion}
\label{sec:discussion}

The improvement mechanism is decoder adaptation \eqref{eq:pes} driven by the
tracking target \eqref{eq:target}: on sustained smooth demand the decoder
acquires a phase-advanced correction that the fixed PD cannot supply, which
is why the largest gains occur on Yaw360 and Full3 and the smallest on
Hardstep. The membrane and gate states are mutually coupled through
\eqref{eq:gate}--\eqref{eq:cond}; the decoder reads gate-derived activity,
but this does not mean $\bm r$ acts alone.

Proposition~\ref{prop:bounded} makes the performance--authority tradeoff
explicit: the attainable deviation scales with $\beta\bm\ell$, so the
conservative hardware caps necessarily bound the achievable improvement.
Principal limitations are single-run mission cells outside Roll, adaptation
remaining active during evaluation, absence of frozen-weight generalization
tests, no armed non-CMG hardware trials, and no completed full-scale
physical yaw360.

\section{Conclusion}
\label{sec:conclusion}

A bounded, baseline-preserving neuromorphic residual with an online
PES-adapted decoder improves quaternion attitude tracking by 17.3--60.6\%
across five missions and two rotor-actuated MuJoCo platforms without rotor
saturation, and transfers qualitatively but with a much smaller margin to a
physical 3DOF~Hover at conservative residual authority. Feasibility and
baseline recovery are structural properties of the residual law rather than
empirical observations, and the learning-rule ablation identifies exactly
the step-command failure mode that these properties remove. Future work:
frozen-weight and held-out evaluation of the annealed rule under the
additive law, repeated physical trials on all
missions, frozen-weight and held-out evaluation, and deployment on
neuromorphic hardware to quantify the energy argument that motivates the
approach.

\section*{Acknowledgment}

% TODO: advisor, laboratory, funding, equipment.
% TODO: ACC 2027 requires transparent disclosure of LLM use. State which
% research, coding, analysis, visualization and writing tasks used which
% tools; human authors retain responsibility for all claims.

\bibliographystyle{IEEEtran}
\bibliography{references}

\end{document}
